% GENERATED FILE. Edit canonical lessons, then run npm run generate:books.
% canonical-source-hash: fnv1a64:debaed0849598809
% canonical-lessons: JA-C148-kaminoke, JA-C148-karu, JA-C148-hige, JA-C148-hayasu, JA-C148-keshou

\chapter{Hair, Beards and Make-up}
\label{ch:ja-hair-beards-and-make-up}
\hlchaptermodality{148}

\begin{chapteropening}
\textbf{By the end of this chapter:} \emph{I can talk about hair, beards and make-up.}

Complete the last lesson of chapter 148: I can talk about hair, beards and make-up.
\end{chapteropening}

\section[\texorpdfstring{\ja{かみのけ} (kaminoke)}{かみのけ (kaminoke)}]{\ja{かみのけ} (kaminoke) — the hair on one's head}

\label{lesson:JA-C148-kaminoke}



\begin{quote}
Before the new one: say the Japanese for hot water, then the Japanese for soap.
\end{quote}

\subsection*{You'll want to know: \ja{かみのけ}}
\textbf{\ja{かみのけ}} — \emph{kaminoke} — \textquotedblleft{}the hair on one's head\textquotedblright{}.

\begin{quote}
\textbf{\ja{いもうとの} \ja{かみのけは} \ja{ながいです。}} — \emph{imōto no kaminoke wa nagai desu} — \textquotedblleft{}My younger sister's hair is long.\textquotedblright{}
\end{quote}

\begin{cousinweb}
\textbf{\ja{かみのけ}} (\emph{kaminoke}) is native Japanese, built from \emph{kami}, \textquotedblleft{}hair\textquotedblright{}, the linking \emph{no}, and \emph{ke}, \textquotedblleft{}hair, fur\textquotedblright{}, all inherited rather than borrowed (JMdict, KANJIDIC2).
\end{cousinweb}

\begin{grammarlens}[title={What you've built}]
One more word for everyday life at A2.
\end{grammarlens}

\subsection*{Your turn}
\begin{itemize}
\raggedright
  \item \emph{Say it:} \emph{kaminoke}
  \item \emph{Say it:} \emph{kaminoke}, once more
  \item \emph{Say it:} \emph{imōto no kaminoke wa nagai desu} — My younger sister's hair is long.
  \item \emph{Recall:} say the Japanese for hot water, then the Japanese for soap, then say \emph{kaminoke} again
\end{itemize}

\begin{tcolorbox}[breakable,colback=teal!4,colframe=teal!35!black,title={Before you move on}]
What does \ja{かみのけ} mean? (\textquotedblleft{}The hair on one's head\textquotedblright{}.) Say it once more.
\end{tcolorbox}
\section[\texorpdfstring{\ja{かる} (karu)}{かる (karu)}]{\ja{かる} (karu) — to trim; to mow}

\label{lesson:JA-C148-karu}



\begin{quote}
Before the new one: say the Japanese for soap, then the Japanese for the hair on one's head.
\end{quote}

\subsection*{You'll want to know: \ja{かる}}
\textbf{\ja{かる}} — \emph{karu} — \textquotedblleft{}to trim; to mow\textquotedblright{}.

\begin{quote}
\textbf{\ja{ちちは} \ja{にわの} \ja{くさを} \ja{かります。}} — \emph{chichi wa niwa no kusa o karimasu} — \textquotedblleft{}My father mows the grass in the garden.\textquotedblright{}
\end{quote}

Polite \textbf{\ja{かります}}, \emph{karimasu}; a godan verb.

\begin{cousinweb}
\textbf{\ja{かる}} (\emph{karu}) is a native Japanese word (wago), inherited rather than borrowed from Chinese; its kanji was attached later as a kun'yomi, and its older origin is not settled (JMdict, KANJIDIC2).
\end{cousinweb}

\begin{grammarlens}[title={What you've built}]
One more word for everyday life at A2.
\end{grammarlens}

\subsection*{Your turn}
\begin{itemize}
\raggedright
  \item \emph{Say it:} \emph{karu}
  \item \emph{Say it:} \emph{karu}, once more
  \item \emph{Say it:} \emph{chichi wa niwa no kusa o karimasu} — My father mows the grass in the garden.
  \item \emph{Recall:} say the Japanese for soap, then the Japanese for the hair on one's head, then say \emph{karu} again
\end{itemize}

\begin{tcolorbox}[breakable,colback=teal!4,colframe=teal!35!black,title={Before you move on}]
What does \ja{かる} mean? (\textquotedblleft{}To trim; to mow\textquotedblright{}.) Say it once more.
\end{tcolorbox}
\section[\texorpdfstring{\ja{ひげ} (hige)}{ひげ (hige)}]{\ja{ひげ} (hige) — a beard; a moustache}

\label{lesson:JA-C148-hige}



\begin{quote}
Before the new one: say the Japanese for the hair on one's head, then the Japanese for to trim.
\end{quote}

\subsection*{You'll want to know: \ja{ひげ}}
\textbf{\ja{ひげ}} — \emph{hige} — \textquotedblleft{}a beard; a moustache\textquotedblright{}.

\begin{quote}
\textbf{\ja{おじいさんは} \ja{ひげが} \ja{ながいです。}} — \emph{ojiisan wa hige ga nagai desu} — \textquotedblleft{}My grandfather has a long beard.\textquotedblright{}
\end{quote}

\begin{cousinweb}
\textbf{\ja{ひげ}} (\emph{hige}) is a native Japanese word (wago), inherited rather than borrowed from Chinese, and its older origin is not settled. It is usually written in kana; the kanji JMdict lists for it were attached later (JMdict).
\end{cousinweb}

\begin{grammarlens}[title={What you've built}]
One more word for everyday life at A2.
\end{grammarlens}

\subsection*{Your turn}
\begin{itemize}
\raggedright
  \item \emph{Say it:} \emph{hige}
  \item \emph{Say it:} \emph{hige}, once more
  \item \emph{Say it:} \emph{ojiisan wa hige ga nagai desu} — My grandfather has a long beard.
  \item \emph{Recall:} say the Japanese for the hair on one's head, then the Japanese for to trim, then say \emph{hige} again
\end{itemize}

\begin{tcolorbox}[breakable,colback=teal!4,colframe=teal!35!black,title={Before you move on}]
What does \ja{ひげ} mean? (\textquotedblleft{}A beard; a moustache\textquotedblright{}.) Say it once more.
\end{tcolorbox}
\section[\texorpdfstring{\ja{はやす} (hayasu)}{はやす (hayasu)}]{\ja{はやす} (hayasu) — to grow a beard; to let grow}

\label{lesson:JA-C148-hayasu}



\begin{quote}
Before the new one: say the Japanese for to trim, then the Japanese for a beard.
\end{quote}

\subsection*{You'll want to know: \ja{はやす}}
\textbf{\ja{はやす}} — \emph{hayasu} — \textquotedblleft{}to grow a beard; to let grow\textquotedblright{}.

\begin{quote}
\textbf{\ja{ふゆは} \ja{ひげを} \ja{はやします。}} — \emph{fuyu wa hige o hayashimasu} — \textquotedblleft{}In winter I grow a beard.\textquotedblright{}
\end{quote}

Polite \textbf{\ja{はやします}}, \emph{hayashimasu}; a godan verb.

\begin{cousinweb}
\textbf{\ja{はやす}} (\emph{hayasu}) is a native Japanese word (wago), inherited rather than borrowed from Chinese; its kanji was attached later as a kun'yomi, and its older origin is not settled (JMdict, KANJIDIC2).
\end{cousinweb}

\begin{grammarlens}[title={What you've built}]
One more word for everyday life at A2.
\end{grammarlens}

\subsection*{Your turn}
\begin{itemize}
\raggedright
  \item \emph{Say it:} \emph{hayasu}
  \item \emph{Say it:} \emph{hayasu}, once more
  \item \emph{Say it:} \emph{fuyu wa hige o hayashimasu} — In winter I grow a beard.
  \item \emph{Recall:} say the Japanese for to trim, then the Japanese for a beard, then say \emph{hayasu} again
\end{itemize}

\begin{tcolorbox}[breakable,colback=teal!4,colframe=teal!35!black,title={Before you move on}]
What does \ja{はやす} mean? (\textquotedblleft{}To grow a beard; to let grow\textquotedblright{}.) Say it once more.
\end{tcolorbox}
\section[\texorpdfstring{\ja{けしょう} (keshō)}{けしょう (keshō)}]{\ja{けしょう} (keshō) — make-up}

\label{lesson:JA-C148-keshou}



\begin{quote}
Before the new one: say the Japanese for a beard, then the Japanese for to grow a beard.
\end{quote}

\subsection*{You'll want to know: \ja{けしょう}}
\textbf{\ja{けしょう}} — \emph{keshō} — \textquotedblleft{}make-up\textquotedblright{}.

\begin{quote}
\textbf{\ja{あねは} \ja{かがみで} \ja{けしょうを} \ja{なおします。}} — \emph{ane wa kagami de keshō o naoshimasu} — \textquotedblleft{}My elder sister fixes her make-up in the mirror.\textquotedblright{}
\end{quote}

\begin{cousinweb}
\textbf{\ja{けしょう}} (\emph{keshō}) is Sino-Japanese (kango): two Chinese characters read with their on'yomi, readings borrowed from Chinese. They are \emph{ke}, \textquotedblleft{}change, take a form\textquotedblright{} (Middle Chinese \emph{xwaeH}) and \emph{shō}, \textquotedblleft{}cosmetics, adorn\textquotedblright{}, whose Middle Chinese reading the sources here do not give (KANJIDIC2; Middle Chinese in Baxter's transcription).
\end{cousinweb}

\begin{grammarlens}[title={What you've built}]
One more word for everyday life at A2.
\end{grammarlens}

\subsection*{Your turn}
\begin{itemize}
\raggedright
  \item \emph{Say it:} \emph{keshō}
  \item \emph{Say it:} \emph{keshō}, once more
  \item \emph{Say it:} \emph{ane wa kagami de keshō o naoshimasu} — My elder sister fixes her make-up in the mirror.
  \item \emph{Recall:} say the Japanese for a beard, then the Japanese for to grow a beard, then say \emph{keshō} again
\end{itemize}

\begin{tcolorbox}[breakable,colback=teal!4,colframe=teal!35!black,title={Before you move on}]
What does \ja{けしょう} mean? (\textquotedblleft{}Make-up\textquotedblright{}.) Say it once more.
\end{tcolorbox}
